% =============================================================================
% macros.tex
% Custom commands and notation for the PhD proposal.
% This file is \input{} into proposal.tex — do NOT add \documentclass here.
%
% HOW TO USE:
%   Define a command once here, then use it throughout all chapter files.
%   To change notation globally, edit only this file.
% =============================================================================


% -----------------------------------------------------------------------------
% 1. GENERAL MATHEMATICS
% Common symbols and operators used across all chapters.
% -----------------------------------------------------------------------------

% Upright 'd' for differentials (e.g. dx, dt) — standard in physics/engineering
\newcommand{\diff}{\mathop{}\!\mathrm{d}}
% Usage: $\int f(x) \diff x$

% Common number sets
\newcommand{\R}{\mathbb{R}}    % Real numbers ℝ
\newcommand{\N}{\mathbb{N}}    % Natural numbers ℕ
\newcommand{\Z}{\mathbb{Z}}    % Integers ℤ

% Norms and absolute value
\newcommand{\norm}[1]{\left\lVert #1 \right\rVert}
\newcommand{\abs}[1]{\left\lvert #1 \right\rvert}
% Usage: $\norm{\mathbf{v}}$   $\abs{x - y}$

% Partial and total derivatives
\newcommand{\pd}[2]{\frac{\partial #1}{\partial #2}}       % \pd{f}{x}
\newcommand{\pdd}[2]{\frac{\partial^2 #1}{\partial #2^2}}  % \pdd{f}{x}
\newcommand{\td}[2]{\frac{\mathrm{d} #1}{\mathrm{d} #2}}  % \td{f}{t}

% Vector and tensor notation
\newcommand{\vect}[1]{\mathbf{#1}}          % bold vector: \vect{u}
\newcommand{\tens}[1]{\boldsymbol{#1}}      % bold greek tensor: \tens{\sigma}
\newcommand{\mat}[1]{\mathbf{#1}}           % matrix: \mat{A}
\newcommand{\transpose}{^{\mathsf{T}}}      % transpose: \mat{A}\transpose


% -----------------------------------------------------------------------------
% 2. CONTINUUM MODEL
% Notation for stress, strain, and field variables.
% -----------------------------------------------------------------------------

\newcommand{\stress}{\tens{\sigma}}            % Cauchy stress tensor σ
\newcommand{\strain}{\tens{\varepsilon}}       % Strain tensor ε
\newcommand{\stiffness}{\mat{C}}               % Stiffness tensor C
\newcommand{\displacement}{\vect{u}}           % Displacement field u
\newcommand{\bodyforce}{\vect{b}}              % Body force vector b
\newcommand{\traction}{\vect{t}}               % Traction vector t

% Nabla / gradient operators
\newcommand{\grad}{\nabla}                     % Gradient ∇
\newcommand{\divg}{\nabla \cdot}               % Divergence ∇·
\newcommand{\lapl}{\nabla^2}                   % Laplacian ∇²

% Domain notation
\newcommand{\domain}{\Omega}                   % Domain Ω
\newcommand{\boundary}{\partial\Omega}         % Boundary ∂Ω


% -----------------------------------------------------------------------------
% 3. MOLECULAR DYNAMICS
% Notation for atomic positions, forces, energies, and potentials.
% -----------------------------------------------------------------------------

\newcommand{\pos}{\vect{r}}                    % Position vector r
\newcommand{\vel}{\vect{v}}                    % Velocity vector v
\newcommand{\force}{\vect{F}}                  % Force vector F
\newcommand{\mass}{m}                          % Atomic mass m
\newcommand{\Natoms}{N}                        % Number of atoms N

% Energies
\newcommand{\Etot}{E_{\mathrm{tot}}}           % Total energy
\newcommand{\Ekin}{E_{\mathrm{kin}}}           % Kinetic energy
\newcommand{\Epot}{E_{\mathrm{pot}}}           % Potential energy
\newcommand{\Eatom}{E_{\mathrm{atom}}}         % Per-atom energy

% Interatomic potential
\newcommand{\phi}{\varphi}                     % Pair potential φ
\newcommand{\cutoff}{r_{\mathrm{c}}}           % Cutoff radius r_c

% Thermodynamic quantities
\newcommand{\Temp}{T}                          % Temperature T
\newcommand{\Press}{P}                         % Pressure P
\newcommand{\kB}{k_{\mathrm{B}}}              % Boltzmann constant k_B

% Mean squared displacement
\newcommand{\MSD}{\langle \abs{\pos(t) - \pos(0)}^2 \rangle}
% Usage: $\MSD$

% Diffusivity
\newcommand{\Diff}{D}                          % Diffusion coefficient D


% -----------------------------------------------------------------------------
% 4. SENSITIVITY ANALYSIS
% Notation for Sobol indices and variance-based methods.
% -----------------------------------------------------------------------------

\newcommand{\model}{\mathcal{M}}               % Model operator M
\newcommand{\inputvec}{\vect{X}}               % Input vector X
\newcommand{\output}{Y}                        % Model output Y

% Variance
\newcommand{\Var}[1]{\mathrm{Var}\!\left(#1\right)}
\newcommand{\Exp}[1]{\mathbb{E}\!\left[#1\right]}
% Usage: $\Var{Y}$   $\Exp{X_i}$

% Sobol sensitivity indices
\newcommand{\Si}{S_i}                          % First-order Sobol index
\newcommand{\Sij}{S_{ij}}                      % Second-order Sobol index
\newcommand{\STi}{S_{T_i}}                     % Total-order Sobol index

% Morris elementary effects
\newcommand{\EE}{\mathrm{EE}}                  % Elementary effect
\newcommand{\mustar}{\mu^*}                    % Mean of |EE| (Morris)
\newcommand{\sigmaEE}{\sigma_{\mathrm{EE}}}    % Std dev of EE


% -----------------------------------------------------------------------------
% 5. WRITING SHORTHANDS
% Abbreviations and text shortcuts for consistent writing style.
% -----------------------------------------------------------------------------
\usepackage{xspace}   % Prevents eating the space after the command

\newcommand{\ie}{i.e.,\xspace}
\newcommand{\eg}{e.g.,\xspace}
\newcommand{\cf}{cf.\xspace}
\newcommand{\etal}{\textit{et al.}\xspace}
\newcommand{\vs}{vs.\xspace}

% Shorthand for referencing equations, figures, tables consistently
% (only needed if NOT using cleveref — if using cleveref, use \cref instead)
% \newcommand{\Eq}[1]{Eq.~(\ref{#1})}
% \newcommand{\Fig}[1]{Fig.~\ref{#1}}
% \newcommand{\Tab}[1]{Table~\ref{#1}}


% -----------------------------------------------------------------------------
% 6. THEOREM-LIKE ENVIRONMENTS
% Numbered environments for definitions, assumptions, and propositions.
% Requires \usepackage{amsthm} in preamble.tex (already included).
% -----------------------------------------------------------------------------
\newtheorem{theorem}{Theorem}[chapter]
\newtheorem{lemma}[theorem]{Lemma}
\newtheorem{corollary}[theorem]{Corollary}
\newtheorem{proposition}[theorem]{Proposition}

\theoremstyle{definition}
\newtheorem{definition}{Definition}[chapter]
\newtheorem{assumption}{Assumption}[chapter]
\newtheorem{example}{Example}[chapter]

\theoremstyle{remark}
\newtheorem*{remark}{Remark}
\newtheorem*{note}{Note}
